\GXNSFBodyText

\noindent\textbf{示例说明}\par
以下仅演示“专业背景—代表性成果—项目任务”的对应关系，姓名、学历、论文、专著、专利和奖励均不得从示例推定。佐佐木希是公开案例对象，不是项目组成员。

\begin{table}[htbp]
  \centering
  \fontsize{10}{14}\selectfont
  \caption{项目组角色与任务分工（占位示例）}
  \begin{tabularx}{\textwidth}{p{0.22\textwidth}X X}
    \toprule
    人员角色 & 专业背景与真实成果填写位置 & 本项目任务 \\
    \midrule
    申请人（请替换） & 【请填写真实学历、职称、代表作及项目经历】 & 科学问题凝练、总体设计、状态模型 \\
    参与者 A（请替换） & 【请填写多语言研究背景及代表性成果】 & 资料核验、实体对齐、编码质控 \\
    参与者 B（请替换） & 【请填写时序建模背景及代表性成果】 & 变点检测、扰动分析、稳健性检验 \\
    \bottomrule
  \end{tabularx}
  \label{tab:gx-team-role}
\end{table}

\noindent\textbf{代表性成果填写提示}\par
正式申请应按模板要求列出论文或专著题目、全部完成人、期刊或出版社、年份、卷期和页码；专利应列出名称、全部完成人、授权地区、类型和授权号；奖励应列出名称、等次、全部完成人和授奖年份，并说明每项成果与本项目研究内容的直接关系。
